\section{Comparison states}\label{sec:comparison-states}

\subsection{Comparison states and translation invariance}\label{subsec:diag-comparison-states}

A comparison state for the chosen cell structure and $O$ is
a law $\nu_0$ obtained by the ordered limits
\begin{equation}\label{eq:diag-comparison-state}
 \mu_{i_m}^{t_k}\longrightarrow\nu_{t_k}
       \quad(m\to\infty)\text{ for each }k,
 \qquad
 \nu_{t_k}\longrightarrow\nu_0\quad(k\to\infty),
 \qquad t_k>0,\quad t_k\to0.
\end{equation}
All convergences here are weak convergences of bosonic laws on
$\Omega$, and the volume subsequence is common to all strengths.
Given a prescribed subsequence of volumes, a further subsequence
produces such a state. To see this, take strengths $t_k=1/k$,
extract nested convergent volume subsequences by the tightness
proved above, and choose their diagonal subsequence. The limiting
laws $\nu_{t_k}$ are themselves uniformly tight: applying
\eqref{eq:diag-limit-domination} to $\beta=\pm1_\sigma$ and
using $t_k\le1$ bounds their exponential coordinate moments
by a common constant. The same product-of-intervals argument
gives a weakly convergent strength subsequence. This proves
existence of \eqref{eq:diag-comparison-state} subordinate to
any prescribed volume subsequence.

Every comparison state has Gaussian domination
\begin{equation}\label{eq:diag-comparison-domination}
 \nu_0(e^{\langle\beta,F\rangle})\le e^{\|\beta\|^2/2},
                         \qquad\beta\text{ of finite support}.
\end{equation}
This follows from \eqref{eq:diag-limit-domination} by letting
$t_k\to0$ and truncating the exponential in the second weak
limit. Wall invariance and wall reflection positivity also
pass through both limits. For each local positive-side test,
its support is contained in the positive half-period slab
once $i_m$ is large, so \eqref{eq:diag-tilted-positive-form}
applies. The continuity and $L^2$ extension described in
\cref{subsec:diag-wall-law} give positivity for bounded Borel
local tests. In particular $\nu_0$ is invariant and reflection
positive at $(x_1-x_2,c_0)$.

The comparison with a prescribed cubic limit will be proved in
\cref{thm:comparison-any-side}, using nearby divisible sides.

We now obtain the extra invariance required to use
\cref{thm:comparison-any-side} for observables at arbitrary
positions. Fix a coordinate link reflection $\theta_c$
and a curvature character $P=\varphi_\beta$. Its polynomial
representation is $e^{i\langle d^*\beta,\widehat A\rangle}$;
on every sufficiently large torus this equals
$e^{i\langle\beta,\widetilde F\rangle}$ on the local lift.
Apply \Cref{thm:diag-polynomial-reduction} to $P$
and its reflection, and then
\Cref{prop:diag-two-layer-principle} to its two-layer
reduction. The finite errors are at most two $O(t)$ terms
plus a finite sum of the terms in
\eqref{eq:diag-two-layer-finite-error}. Along the common
volume subsequence in \eqref{eq:diag-comparison-state},
take volume to infinity, then $t_k\to0$ with the regularization
fixed, and finally remove the regularization. Only the
character expectations are passed to weak limits. The
conclusion is
\begin{equation}\label{eq:diag-character-reflection}
 \nu_0(\Theta_c\varphi_\beta)
                           =\overline{\nu_0(\varphi_\beta)}.
\end{equation}
There is no need to extract limiting fermion expectations
in this argument.

For a bosonic character
$\Theta_c\varphi_\beta=\varphi_{-\theta_c^*\beta}$, while
positivity of the probability law gives
$\overline{\nu_0(\varphi_\beta)}=\nu_0(\varphi_{-\beta})$.
Thus \eqref{eq:diag-character-reflection}, with $-\beta$
in place of $\beta$, is invariance under the tensor pullback
$F\mapsto\theta_c^*F$. Characters determine all finite
plaquette marginals by uniqueness of their Fourier
transforms \cite[Theorem~6.2]{Potzelberger}; these marginals determine the law because
cylinder sets generate the Borel sigma algebra of $\Omega$.
Consequently every comparison state is invariant under
all coordinate link reflections, in all three directions.
The product of reflections at $c$ and $c+1$ is translation
by $2e_\mu$, so
\begin{equation}\label{eq:diag-even-translation}
 (\tau_v^*)_*\nu_0=\nu_0,\qquad v\in2\Z^3.
\end{equation}
We need only these even translations. The original wall
periods alone would not have permitted the next placement.

